\section{Fluid and elastic body}\label{sec_fl_elastic}

In this Section, we will build on the analysis of \cref{sec_fl_rigid} and put \ac{bflu} into interaction with a more complex physical object---an \ac{ela}. Only the main results will be presented here; their derivation follows the lines of  \cref{sec_fl_rigid}.

The \ac{ref} and \ac{cur} configurations are shown in \cref{fig_ref_curr_ela}. Here, the \ac{ela} boundary $\pomcircouteq$ may deform, but the inner \ac{ela} boundary $\pomcircineq$ is pinned to the $\xc^1$, $\xc^2$ plane.



\begin{figure*}
        \centering
        \includegraphics[width=\textwidth]{figures/figure_7/figure_7.pdf}
        \caption{
                \label{fig_ref_curr_ela}
                \Acl{ref} and \acl{cur} configuration for the interaction between a  \acl{bflu} and an \acl{ela}. The figures follows the same notation as \cref{fig_ref_curr_rigid}.
        }
\end{figure*}

%M revise to embed in text
The equations of motion are reported in \cref{app_sec_fl_elastic_equations_of_motion}.



The system dynamics is obtained as follows: given $\bvfr^{n-1}$, $\sigmafr^{n-3/2}$, $\buela^{n-1}$, $\bueladot^{n-1}$, $\budomdot^{n-1}$ and $\budomdot^{n-1}$  from the preceding step, we
\begin{itemize}
        \titleditem{Update \ac{ela}}{ \label{fl_ela_step_1}
                Solve for $\buela^n$ and $\bueladot^n$ the discrete version of \cref{eq_fl_elastic_ela_current}
                \begin{align}
                        \label{eq_fl_elastic_ela_current_disc_1}
                        \rhoela \frac{\ueladot^{n, \, \alpha} -  \ueladot^{n-1, \, \alpha}}{\deltat} & = \pderr{\Sela_{\alpha \beta}(\buela^n)}{\beta}, \\
                        \label{eq_fl_elastic_ela_current_disc_2}
                        \frac{\buela^n -  \buela^{n-1}}{\deltat}                                     & = \bueladot,
                \end{align}
                with \acp{bc} obtained from \cref{bc_fl_ela_1_current,bc_fl_ela_2_current}:
                \bw
                \begin{align}
                        \label{bc_fl_ela_1_current_disc}
                        \buela^n                                                                                                                                                                                                               & =  0          \text{ on }\pomcircineqr, \\
                        \label{bc_fl_ela_2_current_disc}
                        \epsilon_{\beta \gamma} \, \left. \Sela_{\alpha \beta}(\buela^n)\right|_{\bxr = \bxi{s}}\der{\xi^\gamma}{s}\                                                                                                           & =  \newlinenn
                        \sigmastressr_{\alpha \beta}(\bxi{s};\bvfr^{n-1}, \sigmafr^{n-3/2}, \buela^n) \epsilon_{\beta \gamma} \left. F_{\gamma \delta}(\bu_\ela^{n-1})\right|_{\bxr = \bxi{s}}\der{ \xi^\delta}{s}  \text{ on } \pomcircouteqr & .
                \end{align}
                \ew
                In \cref{bc_fl_ela_2_current_disc}, we have used \cref{eq_def_phi,eq_def_F} to rewrite the derivative of ${\bm \phi}(\bxi{s})$.



        }
        \titleditem{Update \ac{dom}}{ \label{fl_ela_step_2}
                Solve for $\budom^n$ and $\budomdot^n$ the \acp{bvp} given by \cref{eq_domain_fluid_rigid_disc,eq_domain_fluid_rigid_disc_dot}  with \acp{bc} \crefs{bc_ela_1_disc}, \crefs{bc_ela_1_disc_dot} and
                \begin{align}
                        \label{bc_ela_3_disc}
                        \budom^n    & = \buela^n \text{ on } \pomcircouteqr,    \\
                        \label{bc_ela_3_dot_disc}
                        \budomdot^n & = \bueladot^n \text{ on } \pomcircouteqr.
                \end{align}
        }


        \titleditem{Update \ac{bflu}}{
                \label{fl_ela_step_3}
                Solve the \acp{bvp} given by  \crefs{eq_fl_rigid_fl_1_reference_disc_aux,eq_poiss_like_phi} and \acp{bc}
                \crefs{bc_fl_rigid_1_reference_disc_aux,bc_fl_rigid_2_reference_disc_aux,bc_fl_rigid_3_reference_disc_aux,bc_fl_rigid_5_reference_disc_aux,bc_phi_1,bc_phi_2} and
                \be
                \label{bc_fl_ela_4_current_disc}
                \bvbarf = \bueladot^n  \text{ on }  \pomcircouteqr,
                \ee
                which replaces \cref{bc_fl_rigid_5_reference_disc_aux} for a rigid body, and solve \cref{eq_phi_1}.


        }\end{itemize}


An example of the coupled dynamics of \ac{bflu} and \ac{ela} is shown in \cref{fig_fluid_ela}, where we show the interaction of a \ac{bflu} with a hydrogel (\ac{ela}).
Here, $\pomcircouteqr$ is an ellipse with semi-axes $a = 0.2\, \met$, $b = 0.1\, \met$ and center located at  $\bxr = (0.4 \, \met, 0.2\, \met)$, and $\pomcircineqr$ has radius $r = 0.0125 \, \met$ and center located at $\bxr$.
The channel dimensions and \ac{bflu} parameters have been taken from the \ac{feat2d} DFG 2D-3 benchmark example for a fluid flow around a cylinder \cite{blumFEAT2DFiniteElement1995,schaferBenchmarkComputationsLaminar1996}.
The \ac{ela} density is $\rhoela = 10^3 \,\kg/\met^2$   \cite{xuConcurrentStiffeningSoftening2023}. The storage modulus, which approximates the shear modulus $\mu$ \cite{REHAGE1974161} is $ \sim  10 \, \kg/\sec^2$ \cite{lichtSynthetic3DPEGAnisogel2019}.
Given a \ac{ela} Poisson ratio \cite{landauTheoryElasticity1986} of $\nu \sim 0.49$ \cite{chippadaSimultaneousDeterminationYoungs2010}, we obtain $K \sim 5 \times 10^2 \, \kg/\sec^2$ from the relation $K = 2 \mu(1+\nu)/[3(2\nu-1)]$ \cite{landauTheoryElasticity1986}.
The \ac{bflu} inflow velocity profile is  a Poiseuille flow, \cref{eq_inflow_profile}, with $ v_0 = 1\, \met / \sec$. Finally, the \ac{dom} mesh stiffness exponent has been chosen to be $\zeta = 3$ \cite{kamenskyLectureNotesMAE}.


\begin{figure*}
        \centering
        \includegraphics[width=\textwidth]{figures/figure_10/figure_10.pdf}
        \caption{
                \label{fig_fluid_ela}
                Interaction between a \acl{bflu} and an \acf{ela} (a hydrogel). The notation is the same as in \cref{fig_fluid_rigid}, and the \acl{ela} is depicted as a red mesh. Here, the Reynolds number is $R \sim 2 \times 10^2$. The \ac{ela} deformation is not symmetric under the mirror symmetry $\xc^2 \rightarrow h-\xc^2$, where $h$ is the height of the \ac{bflu} domain, because the center of \ac{ela} is not located at $\xc^2 = h/2$.
        }
\end{figure*}
